\documentclass[11pt]{article}

\usepackage[T1]{fontenc}
\usepackage{amsmath,amssymb,mathtools}
\usepackage[a4paper,margin=28mm]{geometry}
\usepackage{microtype}
\usepackage[hidelinks]{hyperref}

\newcommand{\dd}{\mathop{}\!\mathrm{d}}
\newcommand{\R}{\mathbb{R}}
\newcommand{\Q}{\mathcal{Q}}
\newcommand{\I}{\mathcal{I}}
\newcommand{\D}{\mathcal{D}}
\newcommand{\be}{\boldsymbol{e}}
\newcommand{\bj}{\boldsymbol{j}}
\newcommand{\bi}{\boldsymbol{i}}
\newcommand{\bt}{\boldsymbol{t}}
\newcommand{\bx}{\boldsymbol{x}}
\newcommand{\balpha}{\boldsymbol{\alpha}}
\newcommand{\bbeta}{\boldsymbol{\beta}}

\title{Quintic Hermite Interpolation\\
  \large Self-contained mathematical notes for one and several dimensions}
\author{}
\date{}

\begin{document}
\maketitle

\begin{abstract}
These notes derive the quintic Hermite interpolant determined by values and
first two derivatives at the ends of an interval.  They then derive the
six-point, value-only variant obtained when those derivatives are replaced by
fourth-order centered differences, and extend both constructions by tensor
products.  Particular care is given to orientation signs, dimensional
scaling, continuity, polynomial reproduction, truncation error, and the
difference between the exact-jet construction of Timmes and Swesty
\cite{timmes-swesty} and the uniform-grid construction described for
\texttt{SPHINCS\_BSSN} by Rosswog and Diener \cite{rosswog-diener}.  Explicit
formulas in the appendix are generated and checked with exact arithmetic by
the accompanying SymPy script.
\end{abstract}

\tableofcontents

\section{Conventions and the two interpolation problems}
\label{sec:conventions}

Let
\[
  x_i<x_{i+1},\qquad h=x_{i+1}-x_i>0,
  \qquad t=\frac{x-x_i}{h}\in[0,1].
\]
Primes on a function denote physical derivatives with respect to $x$, while
primes on a polynomial in $t$, such as $w_j'(t)$, denote derivatives with
respect to $t$.  Thus
\begin{equation}
  \frac{\dd^q}{\dd x^q}=h^{-q}\frac{\dd^q}{\dd t^q}.
  \label{eq:physical-scaling}
\end{equation}

Two related constructions must be kept distinct.
\begin{enumerate}
  \item The \emph{exact-jet interpolant} takes
  $f,f',f''$ at both endpoints as data.  It is the unique polynomial of
  degree at most five satisfying all six conditions.  Its value error is
  $O(h^6)$ for a smooth function.
  \item The \emph{uniform-grid, value-only interpolant} first approximates
  the endpoint derivatives by centered differences and substitutes those
  estimates into the same polynomial.  Each cell polynomial is still
  quintic, but the resulting interpolation operator is generically only
  fifth-order accurate in value.
\end{enumerate}
The first is the mathematical construction used by Timmes and Swesty; the
second is the construction used in the cited \texttt{SPHINCS\_BSSN}
description.  Calling either method ``fifth order'' can therefore refer to
polynomial degree or to the convergence order of the value-only operator; the
meaning should always be stated.

\section{The exact one-dimensional interpolant}
\label{sec:exact-1d}

\subsection{Conditions, monomial form, and uniqueness}

Define the dimensionless endpoint jets
\begin{equation}
  F_{e,r}=h^r f^{(r)}(x_{i+e}),
  \qquad e\in\{0,1\},\quad r\in\{0,1,2\}.
  \label{eq:normalized-jets}
\end{equation}
For $H(t)=\sum_{k=0}^{5}c_k t^k$, the left conditions immediately give
$c_0=F_{0,0}$, $c_1=F_{0,1}$, and $c_2=F_{0,2}/2$.  The right conditions
then give the $3\times3$ system
\begin{equation}
  \begin{pmatrix}
    1&1&1\\ 3&4&5\\ 6&12&20
  \end{pmatrix}
  \begin{pmatrix}c_3\\c_4\\c_5\end{pmatrix}
  =
  \begin{pmatrix}
    F_{1,0}-c_0-c_1-c_2\\
    F_{1,1}-c_1-2c_2\\
    F_{1,2}-2c_2
  \end{pmatrix}.
  \label{eq:monomial-system}
\end{equation}
Its explicit solution is recorded in \hyperref[app:formulas]{Appendix~A}.
This also proves existence.  For uniqueness, the difference of two solutions
has zeros of multiplicity three at both $t=0$ and $t=1$.  A nonzero such
polynomial is divisible by $t^3(t-1)^3$, which has degree six, and hence
cannot have degree at most five.

\subsection{Cardinal basis and the orientation sign}

It is convenient to begin with three left-end basis functions,
\begin{align}
  \psi_0(t)&=1-10t^3+15t^4-6t^5,\nonumber\\
  \psi_1(t)&=t-6t^3+8t^4-3t^5,\label{eq:psi-basis}\\
  \psi_2(t)&=\tfrac12\bigl(t^2-3t^3+3t^4-t^5\bigr).\nonumber
\end{align}
The basis at either endpoint is most safely defined in one formula:
\begin{equation}
  B_{e,r}(t)=(-1)^{er}\,
  \psi_r\!\left(e+(1-2e)t\right),
  \qquad e\in\{0,1\},\quad r\in\{0,1,2\}.
  \label{eq:oriented-basis}
\end{equation}
The factor $(-1)^{er}$ compensates for differentiating the reflected
argument at the right endpoint.  In particular,
\(B_{1,1}(t)=-\psi_1(1-t)\), whereas
\(B_{1,0}(t)=\psi_0(1-t)\) and
\(B_{1,2}(t)=\psi_2(1-t)\).  Direct differentiation gives the cardinal
conditions
\begin{equation}
  \left.\frac{\dd^s B_{e,r}}{\dd t^s}\right|_{t=q}
  =\delta_{eq}\delta_{rs},
  \qquad q,e\in\{0,1\},\quad r,s\in\{0,1,2\}.
  \label{eq:cardinality}
\end{equation}
Consequently the interpolant is
\begin{equation}
  H_i[f](x)=
  \sum_{e=0}^{1}\sum_{r=0}^{2}
  h^r f^{(r)}(x_{i+e})B_{e,r}(t).
  \label{eq:exact-hermite}
\end{equation}
The powers of $h$ are essential: together with
\eqref{eq:physical-scaling} and \eqref{eq:cardinality}, they give
\(H_i^{(r)}(x_{i+e})=f^{(r)}(x_{i+e})\) for $r=0,1,2$.  Expanded basis
derivatives through second order are given in
\hyperref[app:formulas]{Appendix~A}.

\subsection{Remainder, reproduction, and
  \texorpdfstring{$C^2$}{C2} assembly}

The standard repeated-node Hermite remainder is, for each interior $x$,
\begin{equation}
  f(x)-H_i[f](x)
  =\frac{f^{(6)}(\xi_x)}{6!}
    (x-x_i)^3(x-x_{i+1})^3,
  \qquad \xi_x\in(x_i,x_{i+1}),
  \label{eq:exact-remainder}
\end{equation}
under the usual smoothness assumptions.  Equivalently, its magnitude is
$O(h^6)$ at a fixed normalized position.  Every polynomial of degree at most
five is reproduced exactly.

If one polynomial is built on every interval and the same value, first
derivative, and second derivative are assigned to a shared node, the two
adjacent polynomials have identical derivatives through order two there.
The piecewise interpolant is therefore $C^2$.  These assigned jets may be
exact data or consistent approximations; continuity is an algebraic property
of sharing them, not an accuracy statement about them.

\section{Uniform-grid closure using values only}
\label{sec:value-only-1d}

\subsection{Centered derivative formulas from moments}

Now suppose $x_i=x_0+ih$ is a uniform grid and only $f_i=f(x_i)$ is
available.  Write a five-point approximation in the form
\[
  \D_r f_i=h^{-r}\sum_{k=-2}^{2}a_k^{(r)}f_{i+k}.
\]
Matching the Taylor series through degree four imposes
\begin{equation}
  \sum_{k=-2}^{2}a_k^{(r)}k^m=r!\,\delta_{mr},
  \qquad m=0,\ldots,4.
  \label{eq:fd-moments}
\end{equation}
Solving these exact moment equations gives
\begin{align}
  \D_1 f_i
  &=\frac{f_{i-2}-8f_{i-1}+8f_{i+1}-f_{i+2}}{12h},
  \label{eq:fd-first}\\
  \D_2 f_i
  &=\frac{-f_{i-2}+16f_{i-1}-30f_i+16f_{i+1}-f_{i+2}}{12h^2}.
  \label{eq:fd-second}
\end{align}
Their leading defects are
\begin{align}
  \D_1 f_i&=f'_i-\frac{h^4}{30}f_i^{(5)}+O(h^6),
  \label{eq:fd-first-error}\\
  \D_2 f_i&=f''_i-\frac{h^4}{90}f_i^{(6)}+O(h^6).
  \label{eq:fd-second-error}
\end{align}
Thus both derivative estimates are fourth-order accurate.  The cancellation
of the odd moments in \eqref{eq:fd-second} explains why matching only through
degree four still leaves an $O(h^4)$, rather than $O(h^3)$, error.

\subsection{Collection into six grid weights}

Let $a_k=a_k^{(1)}$, $b_k=a_k^{(2)}$, and define them to be zero outside
$-2\leq k\leq2$.  In increasing index order,
\begin{equation}
  (a_{-2},\ldots,a_2)=\frac1{12}(1,-8,0,8,-1),\qquad
  (b_{-2},\ldots,b_2)=\frac1{12}(-1,16,-30,16,-1).
  \label{eq:fd-coefficients}
\end{equation}
Substituting \(\D_1 f_i,\D_2 f_i\) and their shifted counterparts at
$i+1$ into \eqref{eq:exact-hermite}, then collecting values, gives
\begin{equation}
  \I_h f(x_i+th)=\sum_{j=-2}^{3}w_j(t)f_{i+j},
  \label{eq:six-point-interpolant}
\end{equation}
where
\begin{equation}
  \boxed{
  w_j=\delta_{j0}B_{0,0}+\delta_{j1}B_{1,0}
      +a_j B_{0,1}+a_{j-1}B_{1,1}
      +b_j B_{0,2}+b_{j-1}B_{1,2}.}
  \label{eq:collected-weight}
\end{equation}
All basis functions in \eqref{eq:collected-weight} are evaluated at $t$.
The six explicit polynomials and their first two derivatives are given in
\hyperref[app:formulas]{Appendix~A}.

The endpoint values and derivatives of the weights recover precisely the
data and stencils used to build them:
\begin{align}
  w_j(0)&=\delta_{j0},& w_j(1)&=\delta_{j1},\nonumber\\
  w_j'(0)&=a_j,& w_j'(1)&=a_{j-1},\label{eq:weight-endpoints}\\
  w_j''(0)&=b_j,& w_j''(1)&=b_{j-1}.\nonumber
\end{align}
After applying the factors $h^{-1}$ and $h^{-2}$, these are
\eqref{eq:fd-first} and \eqref{eq:fd-second}.  A neighboring cell uses the
same nodal stencils at its common endpoint.  Hence the value-only piecewise
interpolant is exactly $C^2$ at interior grid points, even though its nodal
derivatives are approximations.

\subsection{Symmetry, moments, and truncation error}

The weights obey
\begin{align}
  \sum_{j=-2}^{3}w_j(t)&=1,
  \label{eq:partition-unity}\\
  w_j(t)&=w_{1-j}(1-t),
  \label{eq:weight-reflection}\\
  \sum_{j=-2}^{3}j^m w_j(t)&=t^m,
  \qquad m=0,1,2,3,4.
  \label{eq:weight-moments}
\end{align}
Thus constants are preserved, reflection about the cell midpoint merely
reverses the stencil, and all polynomials through degree four are reproduced.
The first failed moment is
\begin{equation}
  \sum_{j=-2}^{3}j^5w_j(t)-t^5
  =4t(t-1)(2t-1)(3t^2-3t-1).
  \label{eq:fifth-moment-defect}
\end{equation}
This identity is a useful regression test: although each cell expression has
degree five in $t$, the value-only operator is not generally exact on a
quintic function.

Let $E_h=\I_h f-f$, with both terms evaluated at $x_i+th$, and set
\begin{equation}
  A(t)=\frac{t(t-1)(2t-1)(3t^2-3t-1)}{30}.
  \label{eq:A-def}
\end{equation}
Taylor expansion about $x_i$, using \eqref{eq:fifth-moment-defect}, yields
\begin{align}
  E_h&=h^5A(t)f^{(5)}(x_i)+O(h^6),
  \label{eq:value-error}\\
  \frac{\dd E_h}{\dd x}
  &=\frac{h^4}{30}
    (30t^4-60t^3+30t^2-1)f^{(5)}(x_i)+O(h^5),
  \label{eq:first-derivative-error}\\
  \frac{\dd^2E_h}{\dd x^2}
  &=2h^3t(t-1)(2t-1)f^{(5)}(x_i)+O(h^4).
  \label{eq:second-derivative-error}
\end{align}
The generic accuracies of the interpolated value, first derivative, and
second derivative are therefore fifth, fourth, and third order,
respectively.  At a grid endpoint the leading term in
\eqref{eq:second-derivative-error} vanishes and the fourth-order stencil
accuracy \eqref{eq:fd-second-error} is recovered.

At the midpoint the weights reduce to the symmetric vector
\begin{equation}
  \bigl(w_{-2},w_{-1},w_0,w_1,w_2,w_3\bigr)_{t=1/2}
  =\frac1{256}(3,-25,150,150,-25,3).
  \label{eq:midpoint-weights}
\end{equation}
The fifth-degree defect vanishes there, so the value is superconvergent:
\begin{equation}
  E_h\big|_{t=1/2}
  =\frac{5}{1024}h^6f^{(6)}(x_{i+1/2})+O(h^7).
  \label{eq:midpoint-value-error}
\end{equation}
The first derivative remains fourth-order, with leading coefficient
$7h^4f^{(5)}/240$, while the second derivative improves to fourth order:
\begin{equation}
  \left.\frac{\dd^2E_h}{\dd x^2}\right|_{t=1/2}
  =-\frac{259}{5760}h^4f^{(6)}(x_{i+1/2})+O(h^5).
  \label{eq:midpoint-second-error}
\end{equation}

The weights are not all nonnegative; for example
$w_{-1}(1/2)=w_2(1/2)=-25/256$.  The scheme therefore preserves neither
positivity, monotonicity, nor general bounds, and it can overshoot near sharp
features.  Quintic degree and $C^2$ continuity do not imply shape
preservation.

\section{Exact tensor-product Hermite interpolation}
\label{sec:exact-tensor}

\subsection{General rectangular cell}

Consider a $d$-dimensional rectangular cell.  In coordinate $k$, let its
left grid index be $i_k$, its width be $h_k>0$, and
\[
  t_k=\frac{x_k-x_{k,i_k}}{h_k}\in[0,1].
\]
For a corner \(\be\in\{0,1\}^d\), let \(\bx_{\bi+\be}\) denote its physical
position.  For a componentwise derivative multi-index
\(\balpha\in\{0,1,2\}^d\), use
\[
  |\balpha|=\sum_{k=1}^d\alpha_k,\qquad
  \partial^{\balpha}=\prod_{k=1}^d\partial_{x_k}^{\alpha_k},\qquad
    \boldsymbol{h}^{\balpha}=\prod_{k=1}^d h_k^{\alpha_k}.
\]
The exact tensor-product interpolant is
\begin{equation}
  \boxed{
  \I_{\mathrm H}f(\bx)=
  \sum_{\be\in\{0,1\}^d}
  \sum_{\balpha\in\{0,1,2\}^d}
  \boldsymbol{h}^{\balpha}
  (\partial^{\balpha}f)(\bx_{\bi+\be})
  \prod_{k=1}^d B_{e_k,\alpha_k}(t_k).}
  \label{eq:exact-tensor}
\end{equation}
It belongs to
\begin{equation}
  \Q_5^{(d)}
  =\operatorname{span}\left\{
    \bt^{\boldsymbol{\nu}}:0\leq\nu_k\leq5
    \text{ for every }k\right\},
  \qquad \dim\Q_5^{(d)}=6^d.
  \label{eq:Q5}
\end{equation}
Here ``quintic in $d$ dimensions'' means degree at most five in \emph{each}
coordinate; it does not mean total degree at most five.

Taking normalized derivatives and using \eqref{eq:cardinality} shows that
\eqref{eq:exact-tensor} reproduces every prescribed corner datum
\(\partial^{\bbeta}f\), \(\bbeta\in\{0,1,2\}^d\).  There are
$2^d3^d=6^d$ such conditions, matching \eqref{eq:Q5}; tensor cardinality
also proves uniqueness.  Because each factor is a polynomial, mixed
derivatives commute identically.

Across a face shared by two cells, the traces of the normal derivatives of
orders zero, one, and two are the same $(d-1)$-dimensional Hermite
polynomials built from the shared face data.  Tangential differentiation of
those identical traces proves continuity of every derivative of total order
at most two.  Thus a consistently shared grid jet produces a global $C^2$
piecewise tensor interpolant even when the rectangular widths vary from cell
to cell.

\subsection{The 36-condition biquintic}

For $d=2$, each of the four corners supplies the nine fields
\begin{equation}
  \left\{\partial_x^r\partial_y^s f:0\leq r,s\leq2\right\}
  =\{f,f_x,f_y,f_{xx},f_{xy},f_{yy},
      f_{xxy},f_{xyy},f_{xxyy}\}.
  \label{eq:nine-fields}
\end{equation}
Equation \eqref{eq:exact-tensor}, with $d=2$, is therefore a compact and
unambiguous representation of all $4\times9=36$ terms of the biquintic.
Expanding those terms is unnecessary and makes index, sign, and spacing
mistakes harder to detect.

The total-order-two jet has only the first six fields in
\eqref{eq:nine-fields}.  The three higher ``twist'' fields are additional
data needed to select a unique full tensor polynomial in \(\Q_5^{(2)}\).
They are not intrinsically fixed merely by demanding $C^2$ agreement at the
vertices.  To obtain $C^2$ face traces, however, whatever exact values or
closure convention is chosen for them must be shared by neighboring cells.
Exact tabulated jets, a consistent zero convention, or the separable
finite-difference construction of Section~\ref{sec:grid-tensor} all close this
freedom; they generally produce different interiors.

\subsection{The 216-condition triquintic}

For $d=3$, each corner supplies $3^3=27$ fields
\begin{equation}
  \left\{\partial_x^r\partial_y^s\partial_z^q f:
  0\leq r,s,q\leq2\right\},
  \label{eq:27-fields}
\end{equation}
including mixed derivatives up to $f_{xxyyzz}$.  The eight corners give
$8\times27=216$ conditions, exactly the dimension of \(\Q_5^{(3)}\).
Again, \eqref{eq:exact-tensor} is the complete 216-term formula; no expanded
list is mathematically useful.

\section{The separable \texorpdfstring{$6^d$}{6d}-point value-only stencil}
\label{sec:grid-tensor}

On a Cartesian grid that is uniform along each coordinate direction (the
spacings $h_k$ may differ between directions), applying the one-dimensional
value-only operator successively gives
\begin{equation}
  \boxed{
  \I_{\mathrm G}f(\bx)=
  \sum_{\bj\in\{-2,-1,0,1,2,3\}^d}
  f_{\bi+\bj}\prod_{k=1}^d w_{j_k}(t_k).}
  \label{eq:grid-tensor}
\end{equation}
This is a $6^d$-sample stencil: 36 values in two dimensions and 216 values
in three dimensions.  Algebraically, it is also what results from
\eqref{eq:exact-tensor} when every componentwise corner derivative is formed
by applying the centered one-dimensional difference operators in the
corresponding directions.  The tensor form fixes all higher twist data
consistently and makes the order of successive directional interpolations
irrelevant.

For any multi-index \(\boldsymbol{\gamma}\) with nonnegative components,
differentiation is especially simple:
\begin{equation}
  \partial^{\boldsymbol{\gamma}}\I_{\mathrm G}f
  =\sum_{\bj}f_{\bi+\bj}
    \prod_{k=1}^d
    h_k^{-\gamma_k}w_{j_k}^{(\gamma_k)}(t_k).
  \label{eq:grid-mixed-derivative}
\end{equation}
In particular, with $a\neq b$,
\begin{align}
  \partial_{x_a}\I_{\mathrm G}f
  &=\sum_{\bj}f_{\bi+\bj}\,
    h_a^{-1}w'_{j_a}(t_a)
    \prod_{k\neq a}w_{j_k}(t_k),
  \label{eq:grid-gradient}\\
  \partial_{x_a}^2\I_{\mathrm G}f
  &=\sum_{\bj}f_{\bi+\bj}\,
    h_a^{-2}w''_{j_a}(t_a)
    \prod_{k\neq a}w_{j_k}(t_k),
  \label{eq:grid-hessian-diagonal}\\
  \partial_{x_a}\partial_{x_b}\I_{\mathrm G}f
  &=\sum_{\bj}f_{\bi+\bj}\,
    (h_a h_b)^{-1}w'_{j_a}(t_a)w'_{j_b}(t_b)
    \prod_{k\neq a,b}w_{j_k}(t_k).
  \label{eq:grid-hessian-offdiagonal}
\end{align}

Partition of unity and reflection symmetry follow by taking products of
\eqref{eq:partition-unity} and \eqref{eq:weight-reflection}.  The operator
reproduces \(\Q_4^{(d)}\), because each one-dimensional factor reproduces
degrees zero through four.  For comparable spacings, the generic accuracy of
a derivative with each \(\gamma_k\leq2\) is governed by the largest
per-coordinate derivative order, not by \(|\boldsymbol{\gamma}|\): values are
fifth order, gradients and off-diagonal Hessian entries are fourth order, and
diagonal Hessian entries are third order.  At the center of every coordinate
interval, values are sixth order and all Hessian entries are fourth order,
while gradients remain fourth order.  Shared centered stencils and
\eqref{eq:weight-endpoints} give $C^2$ agreement across every interior grid
face.

The numerical counts in the exact-jet and value-only descriptions coincide
but refer to different data.  The Timmes--Swesty biquintic uses 36 independent
corner-jet conditions.  Its direct three-dimensional analogue uses 216
corner-jet conditions.  The \texttt{SPHINCS\_BSSN} construction instead uses
36 or 216 neighboring \emph{function values}, from which all needed jets are
implicitly reconstructed by separable centered differences.

\section{Mathematical errata in the source presentations}
\label{sec:errata}

The compact oriented and tensor notation above was chosen in part to expose
typographical errors in the supplied sources.  The corrections below concern
the formulas, not the numerical conclusions of either paper.

\subsection{Rosswog and Diener}

In the displayed one-dimensional Hermite polynomial of
\cite{rosswog-diener}, the right-end first-derivative term is printed with a
plus sign.  It must be
\begin{equation}
  -h f'_{i+1}\psi_1(1-t),
  \label{eq:sphincs-sign-correction}
\end{equation}
as follows from the chain rule or \eqref{eq:oriented-basis}.  This is not a
cosmetic convention: with $h=1$ and $f(x)=x$, the printed plus-sign
expression equals $x+2\psi_1(1-t)$, so it fails even the linear regression
test.  The paper's centered first-derivative numerator also closes the
subscript with a parenthesis after $i+2$; that character should be a brace,
giving $f_{i+2}$, with the coefficients
otherwise exactly those in \eqref{eq:fd-first}.  In the three-dimensional
stencil prose, independent coordinate indices should be used rather than the
same $i$ for all three axes.

\subsection{Timmes and Swesty}

Equation~(16) of \cite{timmes-swesty} has the correct minus sign on the
reflected right first-derivative basis.  The long 36-term equation~(17),
however, contains several transcription errors; the correct formula is the
$d=2$ specialization of \eqref{eq:exact-tensor}.  In particular:
\begin{itemize}
  \item the density and temperature basis directions are interchanged
  systematically relative to the normalized coordinates in equation~(18)
  and to the Appendix routine;
  \item two first--first mixed entries contain an extra density derivative in
  their labels, and should both be $F_{T\rho}$;
  \item the $F_{\rho\rho T}|_{i+1,j}$ entry has the scale
  $(\Delta T)^2(\Delta\rho)^2$, which should be
  \(\Delta T(\Delta\rho)^2\);
  \item the temperature cell endpoints and the denominator defining the
  temperature coordinate require $T_j,T_{j+1}$, not $T_i,T_{i+1}$; and
  equation~(19) should express $F_{T\rho}=F_{\rho T}$, rather than repeat the
  same ordering on both sides.
\end{itemize}
The printed Appendix routine is likewise biquintic, not ``bicubic.''  In its
first $F_{\rho T T}$ contribution, the array name must denote that third
derivative (\texttt{fdtt}), not the fourth derivative
\texttt{fddtt}; the normalized temperature derivative also divides by the
temperature spacing \texttt{dt}, not the undefined \texttt{dti}.

Finally, derivatives of a biquintic do not in general become ``biquartic''
and then ``bicubic'' in both variables.  In multidegree notation,
\begin{equation}
  H_\rho\in\Q_{4,5},\quad H_T\in\Q_{5,4},\quad
  H_{\rho\rho}\in\Q_{3,5},\quad
  H_{\rho T}\in\Q_{4,4},\quad
  H_{TT}\in\Q_{5,3}.
  \label{eq:anisotropic-degrees}
\end{equation}
Terms such as $\rho^2H_\rho$ in thermodynamic identities can change these
degrees further.  The terms ``biquartic'' and ``bicubic'' in the source are
therefore at best loose shorthand, not exact degree classifications.

\section{Verification boundary and assumptions}
\label{sec:scope}

The accompanying \texttt{derive\_hermite.py} uses exact SymPy rationals to
derive the cardinal bases, solve the monomial and finite-difference moment
systems, collect the six grid weights, and generate
\texttt{generated\_formulas.tex}.  Its checks include endpoint cardinality,
equivalence of basis and monomial forms, exact-jet reproduction through
degree five, failure of the source's plus-sign variant on a linear function,
the finite-difference defects, weight moments and symmetry, midpoint
superconvergence, and matching through the second derivative across cells.
Tensor checks cover all 36 and 216 corner conditions, componentwise degree
bounds, axis exchange, mixed-partial commutation, and shared-face
continuity.  Thus the generated appendix is both a readable formula table and
an exact symbolic audit trail.

The mathematical scope is as follows.
\begin{itemize}
  \item Exact-jet formulas allow arbitrary positive rectangular cell widths.
  The centered value-only closure assumes a Cartesian grid uniform along each
  axis and an interior stencil extending two nodes left and three nodes right
  of the active cell.
  \item Functions are assumed smooth enough for every derivative and error
  term stated.  The formulas are scalar but apply componentwise to vector and
  tensor fields.
  \item Boundary closures, ghost-zone policies, floating-point evaluation,
  operation ordering, optimization, limiting, and production implementation
  are deliberately outside the scope of these notes.
\end{itemize}

\appendix
\section{Generated explicit formulas}
\label{app:formulas}

The following formulas are generated deterministically by the symbolic
derivation script.  They are included rather than copied by hand so that the
document and its executable algebra cannot silently diverge.

% Generated by derive_hermite.py; do not edit by hand.
% Every displayed expression follows from an exact SymPy derivation.

\subsection{One-dimensional cardinal bases}

The left-end cardinal polynomials and the oriented endpoint bases are
\begin{align*}
  \psi_{0}(t) &= - \left(t - 1\right)^{3} \left(6 t^{2} + 3 t + 1\right) \\
  \psi_{1}(t) &= - t \left(t - 1\right)^{3} \left(3 t + 1\right) \\
  \psi_{2}(t) &= - \frac{t^{2} \left(t - 1\right)^{3}}{2}
\end{align*}
\[
  B_{e,r}(t)=(-1)^{er}\psi_r\!\left(e+(1-2e)t\right),
  \qquad e\in\{0,1\},\quad r\in\{0,1,2\}.
\]
Its normalized derivatives obey
\[
  B_{e,r}^{(q)}(t)=(-1)^{e(r+q)}
  \psi_r^{(q)}\!\left(e+(1-2e)t\right).
\]

\subsection{Derivatives of the cardinal polynomials}

\begingroup
\renewcommand{\arraystretch}{1.45}
\[\begin{array}{c|c|l}
r & q & \psi_r^{(q)}(t) \\ \hline
0 & 0 & - \left(t - 1\right)^{3} \left(6 t^{2} + 3 t + 1\right) \\
0 & 1 & - 30 t^{2} \left(t - 1\right)^{2} \\
0 & 2 & - 60 t \left(t - 1\right) \left(2 t - 1\right) \\
\hline
1 & 0 & - t \left(t - 1\right)^{3} \left(3 t + 1\right) \\
1 & 1 & - \left(t - 1\right)^{2} \left(3 t - 1\right) \left(5 t + 1\right) \\
1 & 2 & - 12 t \left(t - 1\right) \left(5 t - 3\right) \\
\hline
2 & 0 & - \frac{t^{2} \left(t - 1\right)^{3}}{2} \\
2 & 1 & - \frac{t \left(t - 1\right)^{2} \left(5 t - 2\right)}{2} \\
2 & 2 & - \left(t - 1\right) \left(10 t^{2} - 8 t + 1\right) \\
\end{array}\]
\endgroup

\subsection{Monomial mapping}

With $F_{e,r}=h^r f^{(r)}(x_{i+e})$ and
$H(t)=\sum_{k=0}^{5}c_k t^k$, the coefficients are
\begin{align*}
  c_{0} &= F_{0,0} \\
  c_{1} &= F_{0,1} \\
  c_{2} &= \frac{F_{0,2}}{2} \\
  c_{3} &= - \frac{20 F_{0,0} + 12 F_{0,1} + 3 F_{0,2} - 20 F_{1,0} + 8 F_{1,1} - F_{1,2}}{2} \\
  c_{4} &= \frac{30 F_{0,0} + 16 F_{0,1} + 3 F_{0,2} - 30 F_{1,0} + 14 F_{1,1} - 2 F_{1,2}}{2} \\
  c_{5} &= - \frac{12 F_{0,0} + 6 F_{0,1} + F_{0,2} - 12 F_{1,0} + 6 F_{1,1} - F_{1,2}}{2}
\end{align*}

\subsection{Uniform-grid six-point weights}

The derived fourth-order centered endpoint formulas are
\begin{align*}
  h f'_i &= \frac{8 f_{i+1} - f_{i+2} - 8 f_{i-1} + f_{i-2}}{12} \\
  h^2 f''_i &= - \frac{30 f_{i} - 16 f_{i+1} + f_{i+2} - 16 f_{i-1} + f_{i-2}}{12}
\end{align*}
Collecting them in $H(t)=\sum_{j=-2}^{3}w_j(t)f_{i+j}$ gives
\begin{align*}
  w_{-2}(t) &= - \frac{t \left(t - 1\right)^{3} \left(5 t + 2\right)}{24} \\
  w_{-1}(t) &= \frac{t \left(t - 1\right) \left(25 t^{3} - 39 t^{2} + 16\right)}{24} \\
  w_{0}(t) &= - \frac{\left(t - 1\right) \left(25 t^{4} - 38 t^{3} - 3 t^{2} + 12 t + 12\right)}{12} \\
  w_{1}(t) &= \frac{t \left(25 t^{4} - 62 t^{3} + 33 t^{2} + 8 t + 8\right)}{12} \\
  w_{2}(t) &= - \frac{t \left(t - 1\right) \left(25 t^{3} - 36 t^{2} - 3 t - 2\right)}{24} \\
  w_{3}(t) &= \frac{t^{3} \left(t - 1\right) \left(5 t - 7\right)}{24}
\end{align*}

The explicit first normalized derivatives are
\begin{align*}
  w'_{-2}(t) &= - \frac{\left(t - 1\right)^{2} \left(25 t^{2} - 2 t - 2\right)}{24} \\
  w'_{-1}(t) &= \frac{125 t^{4} - 256 t^{3} + 117 t^{2} + 32 t - 16}{24} \\
  w'_{0}(t) &= - \frac{t \left(125 t^{3} - 252 t^{2} + 105 t + 30\right)}{12} \\
  w'_{1}(t) &= \frac{\left(t - 1\right) \left(125 t^{3} - 123 t^{2} - 24 t - 8\right)}{12} \\
  w'_{2}(t) &= - \frac{125 t^{4} - 244 t^{3} + 99 t^{2} + 2 t + 2}{24} \\
  w'_{3}(t) &= \frac{t^{2} \left(25 t^{2} - 48 t + 21\right)}{24}
\end{align*}

The explicit second normalized derivatives are
\begin{align*}
  w''_{-2}(t) &= - \frac{\left(t - 1\right) \left(50 t^{2} - 28 t - 1\right)}{12} \\
  w''_{-1}(t) &= \frac{250 t^{3} - 384 t^{2} + 117 t + 16}{12} \\
  w''_{0}(t) &= - \frac{250 t^{3} - 378 t^{2} + 105 t + 15}{6} \\
  w''_{1}(t) &= \frac{250 t^{3} - 372 t^{2} + 99 t + 8}{6} \\
  w''_{2}(t) &= - \frac{250 t^{3} - 366 t^{2} + 99 t + 1}{12} \\
  w''_{3}(t) &= \frac{t \left(50 t^{2} - 72 t + 21\right)}{12}
\end{align*}

Physical derivatives follow from
$\mathrm{d}^q H/\mathrm{d}x^q=h^{-q}\sum_j w_j^{(q)}(t)f_{i+j}$.


\begin{thebibliography}{9}

\bibitem{rosswog-diener}
S.~Rosswog and P.~Diener,
``\texttt{SPHINCS\_BSSN}: A general relativistic Smooth Particle
Hydrodynamics code for dynamical spacetimes,''
\emph{Classical and Quantum Gravity} \textbf{38} (2021) 115002,
\href{https://doi.org/10.1088/1361-6382/abee65}{doi:10.1088/1361-6382/abee65}.

\bibitem{timmes-swesty}
F.~X. Timmes and F.~D. Swesty,
``The Accuracy, Consistency, and Speed of an Electron-Positron Equation of
State Based on Table Interpolation of the Helmholtz Free Energy,''
\emph{The Astrophysical Journal Supplement Series} \textbf{126} (2000)
501--516,
\href{https://doi.org/10.1086/313304}{doi:10.1086/313304}.

\end{thebibliography}

\end{document}
